% GENERATED FILE. Edit canonical lessons, then run npm run generate:books.
% canonical-source-hash: fnv1a64:c79d1e87216f466f
% canonical-lessons: RU-C330-smeshivat, RU-C330-vzbivat, RU-C330-venchik, RU-C330-pomeshivat, RU-C330-polovnik

\chapter{Mixing and stirring}
\label{ch:ru-mixing-and-stirring}
\hlchaptermodality{330}

\begin{chapteropening}
\textbf{By the end of this chapter:} \emph{I can describe mixing food: mixing, whisking, a whisk, stirring and a ladle.}

Complete the last lesson of chapter 330: I can describe mixing food: mixing, whisking, a whisk, stirring and a ladle.
\end{chapteropening}

\section[\texorpdfstring{\ru{смешивать} (sméshivat)}{смешивать (sméshivat)}]{\ru{смешивать} (sméshivat) — to mix}

\label{lesson:RU-C330-smeshivat}



\begin{quote}
Before the new one: say the Russian for a lid, then the Russian for to cover.
\end{quote}

\subsection*{You'll want to know: \ru{смешивать}}
\textbf{\ru{смешивать}} (\emph{sméshivat}) — \textquotedblleft{}to mix\textquotedblright{}.

\begin{quote}
\ru{Она} \ru{смешала} \ru{муку} \ru{с} \ru{сахаром}.
\end{quote}

\begin{quote}
\emph{Oná smeshála múku s sákharom.}
\end{quote}

\begin{quote}
She mixed the flour with the sugar.
\end{quote}

Imperfective; perfective \emph{\ru{смешать}} (\emph{smeshát}). It takes \emph{\ru{с}} + instrumental.

\begin{cousinweb}
\emph{\ru{смешивать}} is built on \emph{\ru{с}-} and \emph{\ru{мешать}} \textquotedblleft{}to stir, to mix\textquotedblright{}, a cousin of \emph{\ru{месить}} \textquotedblleft{}to knead\textquotedblright{}, which EtymDB traces to PIE \emph{mey\'{k}-} \textquotedblleft{}to mix\textquotedblright{}.
\end{cousinweb}

\begin{grammarlens}[title={What you've built}]
One more word for telling, explaining and describing at B1.
\end{grammarlens}

\subsection*{Your turn}
\begin{itemize}
\raggedright
  \item \emph{Say it:} \emph{sméshivat}
  \item \emph{Say it:} \emph{sméshivat}, once more
  \item \emph{Say it:} a sentence of your own with \emph{sméshivat}
  \item \emph{Recall:} say the Russian for a lid, then the Russian for to cover, then say \emph{sméshivat} again
\end{itemize}

\begin{tcolorbox}[breakable,colback=teal!4,colframe=teal!35!black,title={Before you move on}]
What does \ru{смешивать} mean? (\textquotedblleft{}To mix\textquotedblright{}.) Say it once more.
\end{tcolorbox}
\section[\texorpdfstring{\ru{взбивать} (vzbivát)}{взбивать (vzbivát)}]{\ru{взбивать} (vzbivát) — to whisk, to beat}

\label{lesson:RU-C330-vzbivat}



\begin{quote}
Before the new one: say the Russian for to cover, then the Russian for to mix.
\end{quote}

\subsection*{You'll want to know: \ru{взбивать}}
\textbf{\ru{взбивать}} (\emph{vzbivát}) — \textquotedblleft{}to whisk, to beat\textquotedblright{}.

\begin{quote}
\ru{Папа} \ru{взбивает} \ru{яйца}.
\end{quote}

\begin{quote}
\emph{Pápa vzbiváyet yáytsa.}
\end{quote}

\begin{quote}
Dad is beating the eggs.
\end{quote}

Imperfective; the perfective is \emph{\ru{взбить}} (\emph{vzbit}), \emph{\ru{взобью}}. Also \emph{\ru{взбить} \ru{подушку}} — to plump up a pillow.

\begin{cousinweb}
\emph{\ru{взбивать}} is built on \emph{\ru{вз}-} \textquotedblleft{}up\textquotedblright{} and \emph{\ru{бить}} \textquotedblleft{}to beat\textquotedblright{}, inherited from Proto-Slavic \emph{biti}, PIE \emph{b\textsuperscript{h}eyH-} (EtymDB).
\end{cousinweb}

\begin{grammarlens}[title={What you've built}]
One more word for telling, explaining and describing at B1.
\end{grammarlens}

\subsection*{Your turn}
\begin{itemize}
\raggedright
  \item \emph{Say it:} \emph{vzbivát}
  \item \emph{Say it:} \emph{vzbivát}, once more
  \item \emph{Say it:} a sentence of your own with \emph{vzbivát}
  \item \emph{Recall:} say the Russian for to cover, then the Russian for to mix, then say \emph{vzbivát} again
\end{itemize}

\begin{tcolorbox}[breakable,colback=teal!4,colframe=teal!35!black,title={Before you move on}]
What does \ru{взбивать} mean? (\textquotedblleft{}To whisk, to beat\textquotedblright{}.) Say it once more.
\end{tcolorbox}
\section[\texorpdfstring{\ru{венчик} (vénchik)}{венчик (vénchik)}]{\ru{венчик} (vénchik) — a whisk}

\label{lesson:RU-C330-venchik}



\begin{quote}
Before the new one: say the Russian for to mix, then the Russian for to whisk.
\end{quote}

\subsection*{You'll want to know: \ru{венчик}}
\textbf{\ru{венчик}} (\emph{vénchik}) — \textquotedblleft{}a whisk\textquotedblright{}.

\begin{quote}
\ru{Венчик} \ru{лежит} \ru{в} \ru{ящике}.
\end{quote}

\begin{quote}
\emph{Vénchik lezhít v yáshchike.}
\end{quote}

\begin{quote}
The whisk is in the drawer.
\end{quote}

Diminutive of \emph{\ru{венец}} \textquotedblleft{}a crown, a wreath\textquotedblright{}; also a halo or a corolla.

\begin{cousinweb}
\emph{\ru{венчик}} is a diminutive of \emph{\ru{венец}} \textquotedblleft{}a crown, a wreath\textquotedblright{} (EtymDB lists \emph{\ru{венец}} with no parent): a whisk is a little ring of wire.
\end{cousinweb}

\begin{grammarlens}[title={What you've built}]
One more word for telling, explaining and describing at B1.
\end{grammarlens}

\subsection*{Your turn}
\begin{itemize}
\raggedright
  \item \emph{Say it:} \emph{vénchik}
  \item \emph{Say it:} \emph{vénchik}, once more
  \item \emph{Say it:} a sentence of your own with \emph{vénchik}
  \item \emph{Recall:} say the Russian for to mix, then the Russian for to whisk, then say \emph{vénchik} again
\end{itemize}

\begin{tcolorbox}[breakable,colback=teal!4,colframe=teal!35!black,title={Before you move on}]
What does \ru{венчик} mean? (\textquotedblleft{}A whisk\textquotedblright{}.) Say it once more.
\end{tcolorbox}
\section[\texorpdfstring{\ru{помешивать} (poméshivat)}{помешивать (poméshivat)}]{\ru{помешивать} (poméshivat) — to stir (now and then)}

\label{lesson:RU-C330-pomeshivat}



\begin{quote}
Before the new one: say the Russian for to whisk, then the Russian for a whisk.
\end{quote}

\subsection*{You'll want to know: \ru{помешивать}}
\textbf{\ru{помешивать}} (\emph{poméshivat}) — \textquotedblleft{}to stir (now and then)\textquotedblright{}.

\begin{quote}
\ru{Суп} \ru{нужно} \ru{помешивать}.
\end{quote}

\begin{quote}
\emph{Sup núzhno poméshivat.}
\end{quote}

\begin{quote}
The soup needs stirring.
\end{quote}

Imperfective only: to stir from time to time.

\begin{cousinweb}
\emph{\ru{помешивать}} is built on \emph{\ru{по}-} and \emph{\ru{мешать}} \textquotedblleft{}to stir\textquotedblright{}, a cousin of \emph{\ru{месить}} \textquotedblleft{}to knead\textquotedblright{}, which EtymDB traces to PIE \emph{mey\'{k}-} \textquotedblleft{}to mix\textquotedblright{}.
\end{cousinweb}

\begin{grammarlens}[title={What you've built}]
One more word for telling, explaining and describing at B1.
\end{grammarlens}

\subsection*{Your turn}
\begin{itemize}
\raggedright
  \item \emph{Say it:} \emph{poméshivat}
  \item \emph{Say it:} \emph{poméshivat}, once more
  \item \emph{Say it:} a sentence of your own with \emph{poméshivat}
  \item \emph{Recall:} say the Russian for to whisk, then the Russian for a whisk, then say \emph{poméshivat} again
\end{itemize}

\begin{tcolorbox}[breakable,colback=teal!4,colframe=teal!35!black,title={Before you move on}]
What does \ru{помешивать} mean? (\textquotedblleft{}To stir (now and then)\textquotedblright{}.) Say it once more.
\end{tcolorbox}
\section[\texorpdfstring{\ru{половник} (polóvnik)}{половник (polóvnik)}]{\ru{половник} (polóvnik) — a ladle}

\label{lesson:RU-C330-polovnik}



\begin{quote}
Before the new one: say the Russian for a whisk, then the Russian for to stir.
\end{quote}

\subsection*{You'll want to know: \ru{половник}}
\textbf{\ru{половник}} (\emph{polóvnik}) — \textquotedblleft{}a ladle\textquotedblright{}.

\begin{quote}
\ru{Бабушка} \ru{налила} \ru{суп} \ru{половником}.
\end{quote}

\begin{quote}
\emph{Bábushka nalilá sup polóvnikom.}
\end{quote}

\begin{quote}
Grandma ladled out the soup.
\end{quote}

Masculine; genitive \emph{\ru{половника}}.

\begin{cousinweb}
The origin of \emph{\ru{половник}} is not settled; EtymDB lists no parent.
\end{cousinweb}

\begin{grammarlens}[title={What you've built}]
One more word for telling, explaining and describing at B1.
\end{grammarlens}

\subsection*{Your turn}
\begin{itemize}
\raggedright
  \item \emph{Say it:} \emph{polóvnik}
  \item \emph{Say it:} \emph{polóvnik}, once more
  \item \emph{Say it:} a sentence of your own with \emph{polóvnik}
  \item \emph{Recall:} say the Russian for a whisk, then the Russian for to stir, then say \emph{polóvnik} again
\end{itemize}

\begin{tcolorbox}[breakable,colback=teal!4,colframe=teal!35!black,title={Before you move on}]
What does \ru{половник} mean? (\textquotedblleft{}A ladle\textquotedblright{}.) Say it once more.
\end{tcolorbox}
